\documentclass[a4paper,12pt]{article}

\usepackage[utf8]{inputenc}
\usepackage[T1]{fontenc}
\usepackage[table]{xcolor}
\usepackage{geometry}
\usepackage{longtable}
\usepackage{enumitem}
\usepackage{hyperref}
\usepackage{listings}
\usepackage{titlesec}
\usepackage{booktabs}
\usepackage{array}

\geometry{margin=2.5cm}

\definecolor{criticalcolor}{HTML}{DC3545}
\definecolor{highcolor}{HTML}{FD7E14}
\definecolor{mediumcolor}{HTML}{FFC107}
\definecolor{lowcolor}{HTML}{28A745}
\definecolor{resolvedcolor}{HTML}{6C757D}

\newcommand{\severitybox}[2]{%
  \colorbox{#1}{\makebox[2.5cm][c]{\textcolor{white}{\textbf{#2}}}}%
}

\lstset{
  basicstyle=\small\ttfamily,
  breaklines=true,
  frame=single,
  numbers=none,
  tabsize=2,
  backgroundcolor=\color{gray!10},
  extendedchars=true,
  inputencoding=utf8,
  literate={—}{---}1 {·}{$\cdot$}1 {≈}{$\approx$}1 {×}{$\times$}1
}

\title{\textbf{Performance Report \#09}\\[4pt]
\large Vulkan Engine — Shader Performance Audit}
\author{Henrique Rocha}
\date{\today}

\begin{document}
\maketitle

\begin{abstract}
Unlike previous reports, which focused on CPU-side host code, synchronization
and resource lifetime, this audit is dedicated exclusively to the
\textbf{GPU shader programs} under \texttt{shaders/}. Every finding below is a
pure \emph{performance} concern: none of them require any change to the
rendered result. All existing features — triplanar mapping, cascaded EVSM
shadows, tessellated displacement, procedural water (refraction, reflection,
caustics, glitter), GPU vegetation and every debug visualization mode — are to
be \textbf{preserved unchanged}. The recommendations only remove redundant work,
hoist invariant computation out of per-fragment scope, and gate
development-only code behind specialization constants. 12 issues are reported
(3 high, 6 medium, 3 low), concentrated in the two heaviest fragment shaders
(\texttt{water.frag}, \texttt{main.frag}) and the shared procedural-noise
library.
\end{abstract}

\tableofcontents
\newpage

\section{Scope \& Methodology}

This audit covers the shader tree only:
\texttt{shaders/*.\{vert,frag,geom,tesc,tese,comp\}} and the shared
\texttt{shaders/includes/*.glsl}. Each shader was read in full and analysed for
per-invocation cost, focusing on the fragment and tessellation stages where
invocation counts are highest. Cost estimates assume the default configuration
(3 cascades, tessellation enabled, procedural water with caustics). The guiding
constraint from \texttt{AGENTS.md} is respected: \emph{functionality must be
kept}; only work that is provably redundant or hoistable is targeted.

\subsection*{Shader Hotspot Overview}

\begin{center}
\begin{tabular}{lcl}
\toprule
\textbf{Shader} & \textbf{Lines} & \textbf{Dominant cost} \\
\midrule
\texttt{water.frag}      & 611 & 4D FBM noise, matrix inverse, caustics \\
\texttt{main.frag}       & 514 & triplanar sampling, env cubemap, debug chain \\
\texttt{water.tese}      & 166 & 5$\times$ FBM displacement per vertex \\
\texttt{vegetation.geom} & 187 & already well-optimized (per-plane wind) \\
\texttt{includes/perlin.glsl} & 131 & \texttt{sin()}-based hash, 8/16 taps \\
\bottomrule
\end{tabular}
\end{center}

\newpage

\section{Issue Register}

\subsection{High Severity}

\begin{enumerate}[leftmargin=*]

\item[\textbf{[SH1]}] \textbf{Per-Fragment 4$\times$4 Matrix Inverse in \texttt{water.frag}}
  \hfill \severitybox{highcolor}{HIGH}

  \begin{tabular}{ll}
    \textbf{File:} & \texttt{shaders/water.frag} \\
    \textbf{Line:} & 97 \\
    \textbf{Category:} & Fragment ALU / Invariant Hoisting \\
  \end{tabular}

  \textbf{Description:}
  Every water fragment recomputes a full inverse of the view-projection matrix:

\begin{lstlisting}
mat4 invViewProj = inverse(ubo.viewProjection);
\end{lstlisting}

  \texttt{inverse(mat4)} expands to roughly 150--200 scalar ops (cofactors +
  reciprocal determinant). This is entirely camera-constant: it is identical for
  every fragment of every water pixel in the frame, yet it is evaluated
  millions of times. The result is used only to reconstruct two world-space
  positions (lines 98--101).

  \textbf{Recommendation:}
  Compute the inverse once on the CPU and upload it in the UBO. The
  \texttt{WaterUBO} used by \texttt{postprocess.frag} already carries an
  \texttt{invViewProjection} field, so the pattern and the CPU value already
  exist — add the same \texttt{mat4} to \texttt{SolidParamsUBO} (or the water
  render UBO) and read it directly. Identical output, one matrix multiply
  instead of a per-fragment inverse.

\item[\textbf{[SH2]}] \textbf{Procedural Noise Explosion in Water Caustics/Refraction}
  \hfill \severitybox{highcolor}{HIGH}

  \begin{tabular}{ll}
    \textbf{File:} & \texttt{shaders/water.frag}, \texttt{includes/water\_noise.glsl} \\
    \textbf{Line:} & 285--296, 413--468 \\
    \textbf{Category:} & Fragment ALU \\
  \end{tabular}

  \textbf{Description:}
  The Perlin caustic path (lines 435--468) evaluates
  \texttt{waterRefractionNoise} up to \textbf{5 times} per fragment (front
  centre reused, plus 2 tangent perturbations, plus back centre and 2 more
  perturbations). Each \texttt{waterRefractionNoise} call itself performs
  \emph{four} \texttt{waterFbmNoise} evaluations (lines 24--33), and each FBM
  runs \texttt{noiseOctaves} iterations of \texttt{perlinNoise4D} — 16
  \texttt{hash44} evaluations per octave. On top of that the shader adds
  specular noise (line 285), two glitter FBM calls (lines 291--294) and the
  refraction noise (line 188). A single deep-water fragment can therefore issue
  \textbf{hundreds} of 4D gradient-noise evaluations.

  \textbf{Recommendation (feature-preserving):}
  \begin{itemize}[noitemsep]
    \item Skip the entire caustic block when \texttt{causticIntensity} $\approx 0$
          (see also SH9) — no visual change where caustics are off.
    \item Reduce the caustic Jacobian to a single-sided estimate blended by
          \texttt{depthInfluence} only when the back face is valid; the flat-water
          case (\texttt{hasValidBackFace == false}) never needs the back-face
          FBM set.
    \item Compute the shared FBM lattice once and derive the perturbed samples
          via analytic gradient rather than re-sampling at $\pm\varepsilon$
          (finite differences already assume smoothness).
  \end{itemize}
  These reduce redundant octave work while keeping the same wave/caustic look.

\item[\textbf{[SH3]}] \textbf{\texttt{sin()}-Based Hash in Noise Library}
  \hfill \severitybox{highcolor}{HIGH}

  \begin{tabular}{ll}
    \textbf{File:} & \texttt{includes/perlin.glsl}, \texttt{includes/voronoi.glsl} \\
    \textbf{Line:} & perlin.glsl 4--17; voronoi.glsl 3--9 \\
    \textbf{Category:} & Fragment/Tessellation ALU \\
  \end{tabular}

  \textbf{Description:}
  All gradient/feature hashing uses the classic
  \texttt{fract(sin(dot(...)) * 43758.5453)} idiom:

\begin{lstlisting}
return -1.0 + 2.0 * fract(sin(p) * 43758.5453123);
\end{lstlisting}

  \texttt{hash33} issues 3 \texttt{sin} calls, \texttt{hash44} issues 4. Since
  \texttt{perlinNoise3D} calls \texttt{hash33} 8 times and
  \texttt{perlinNoise4D} calls \texttt{hash44} 16 times, a single 4D noise
  octave costs 64 transcendental \texttt{sin} evaluations. \texttt{sin} is a
  multi-instruction special-function-unit op and is a well-known bottleneck for
  procedural noise. This hash also has known precision artefacts at large
  coordinates on some vendors.

  \textbf{Recommendation:}
  Replace the \texttt{sin} hash with an integer/bit-mixing hash (e.g.\ a
  Wang/PCG-style \texttt{uint} mix on \texttt{floatBitsToUint(p)}), which
  produces statistically equivalent noise with pure integer ALU. This is a
  drop-in change to \texttt{hash33}/\texttt{hash44}/\texttt{hash3} and benefits
  water, terrain displacement and vegetation simultaneously without altering the
  noise character (only the numeric seed constants change).

\end{enumerate}

\newpage

\subsection{Medium Severity}

\begin{enumerate}[leftmargin=*]

\item[\textbf{[SH4]}] \textbf{Redundant \texttt{screenUV} and Scene-Depth Samples in \texttt{water.frag}}
  \hfill \severitybox{mediumcolor}{MEDIUM}

  \begin{tabular}{ll}
    \textbf{File:} & \texttt{shaders/water.frag} \\
    \textbf{Line:} & 85, 89, 141--142, 182, 249 \\
    \textbf{Category:} & Fragment ALU / Texture Bandwidth \\
  \end{tabular}

  \textbf{Description:}
  The clip-to-screen conversion
  \texttt{(fragPosClip.xy / fragPosClip.w) * 0.5 + 0.5} is recomputed at least
  four times (lines 85, 141, 182, and again in debug at 514/521) producing the
  identical value each time. Worse, \texttt{sceneDepthTex} is sampled three
  times at effectively the same coordinate: line 89 (\texttt{earlyScreenUV0}),
  line 142 (\texttt{earlyScreenUV}) and line 249 (\texttt{screenUV}). These are
  the same texel fetched repeatedly.

  \textbf{Recommendation:}
  Compute \texttt{screenUV} once at the top of \texttt{main()} and reuse it;
  sample \texttt{sceneDepthTex} once into a local and reuse the raw value.
  Bit-identical output, fewer texture fetches and address computations per
  fragment.

\item[\textbf{[SH5]}] \textbf{Environment Cubemap Sampled Even When Reflection Strength Is Zero}
  \hfill \severitybox{mediumcolor}{MEDIUM}

  \begin{tabular}{ll}
    \textbf{File:} & \texttt{shaders/main.frag} \\
    \textbf{Line:} & 216--228 \\
    \textbf{Category:} & Texture Bandwidth \\
  \end{tabular}

  \textbf{Description:}
  The environment reflection block samples the cubemap unconditionally whenever
  \texttt{ubo.materialFlags.x < 0.5} (i.e.\ not during 360 capture):

\begin{lstlisting}
vec3 envColor = texture(environmentMap, normalize(envReflectDir)).rgb;
...
blendedRefStrength = refStrength0*w.x + refStrength1*w.y + refStrength2*w.z;
envReflection = envColor * envBlend * blendedRefStrength;
\end{lstlisting}

  For the overwhelming majority of terrain materials \texttt{reflectionStrength}
  (\texttt{tessLevelParams.z}) is 0, so the cubemap sample and the subsequent
  \texttt{reflect}/\texttt{normalize} are pure waste — the result is multiplied
  by zero and \texttt{mix()} at line 503 collapses to the lit colour.

  \textbf{Recommendation:}
  Compute \texttt{blendedRefStrength} first (it is a cheap dot product on data
  already loaded), and only sample the cubemap when it exceeds a small epsilon.
  Because \texttt{blendedRefStrength} is per-material and largely uniform across
  a draw, the branch is highly coherent. Identical output; skips a cubemap fetch
  on all non-reflective surfaces.

\item[\textbf{[SH6]}] \textbf{Triplanar UV Recomputation and Repeated Sampling}
  \hfill \severitybox{mediumcolor}{MEDIUM}

  \begin{tabular}{ll}
    \textbf{File:} & \texttt{includes/triplanar.glsl}, \texttt{shaders/main.frag} \\
    \textbf{Line:} & triplanar.glsl 22--101 \\
    \textbf{Category:} & Fragment ALU / Texture Bandwidth \\
  \end{tabular}

  \textbf{Description:}
  In the triplanar path, \texttt{computeTriplanarAlbedo},
  \texttt{computeTriplanarNormal}, \texttt{computeTriplanarRoughness} and
  \texttt{computeTriplanarAO} each independently call
  \texttt{computeTriplanarUVs} for the \emph{same} \texttt{fragPosWorld} and
  \texttt{brushIndex}, so the projected UVs (with their branch on
  \texttt{geomN} sign) are recomputed four times per material. Because
  \texttt{main.frag} blends three materials, that is up to \textbf{12}
  UV-recomputations and up to \textbf{36} texture samples (3 projections
  $\times$ 3 materials $\times$ up to 4 maps) per fragment.

  \textbf{Recommendation:}
  Compute the three projected UV sets once per material (a single call
  producing \texttt{uvX/uvY/uvZ}) and pass them into the albedo/normal/
  roughness/AO samplers, or fold the four map fetches into one helper that
  reuses the UVs. The per-projection zero-weight guards (\texttt{triW.x > 0.0})
  already skip unused projections — preserve them. No change to blend result.

\item[\textbf{[SH7]}] \textbf{Debug Visualization Chains Compiled Into Release Shaders}
  \hfill \severitybox{mediumcolor}{MEDIUM}

  \begin{tabular}{ll}
    \textbf{File:} & \texttt{shaders/main.frag}, \texttt{shaders/water.frag} \\
    \textbf{Line:} & main.frag 230--499; water.frag 481--604 \\
    \textbf{Category:} & Shader Size / ICache / ALU \\
  \end{tabular}

  \textbf{Description:}
  \texttt{main.frag} contains \textbf{34} sequential \texttt{if (debugMode == N)}
  branches and \texttt{water.frag} contains \textbf{17} \texttt{if (dbgMode == N)}
  branches. Several of these blocks perform texture samples and even full
  triplanar/noise recomputation (e.g.\ main.frag modes 17--30 call
  \texttt{computeTriplanarNormal}; water.frag mode 32 recomputes
  \texttt{waterWaveDisplacement}). Although \texttt{debugMode} is uniform and the
  branches are coherent, the entire chain is compiled into the production
  pipeline: it bloats the shader, keeps live registers for otherwise-dead paths,
  and every fragment still evaluates the comparison ladder.

  \textbf{Recommendation:}
  Gate the debug blocks behind a \texttt{layout(constant\_id = N) const int}
  specialization constant (or a preprocessor \texttt{\#ifdef DEBUG\_MODES}). In
  release pipelines the compiler then dead-code-eliminates the whole ladder,
  shrinking the shader and freeing registers, while debug builds keep every
  visualization mode exactly as-is.

\item[\textbf{[SH8]}] \textbf{Five Displacement Evaluations Per Water Tessellation Vertex}
  \hfill \severitybox{mediumcolor}{MEDIUM}

  \begin{tabular}{ll}
    \textbf{File:} & \texttt{shaders/water.tese} \\
    \textbf{Line:} & 122--149 \\
    \textbf{Category:} & Tessellation ALU \\
  \end{tabular}

  \textbf{Description:}
  Each tessellated water vertex calls \texttt{waterWaveDisplacement} five times:
  once for the height (line 122) and four times for the central-difference
  normal (lines 141--144, at $\pm\varepsilon T$ and $\pm\varepsilon B$). Each
  call is a two-FBM evaluation. With high tessellation factors the vertex count
  is large, making this the dominant tessellation cost. The same 4$\times$
  finite-difference pattern is repeated in \texttt{water.frag} (lines 162--169)
  for the no-tessellation fallback.

  \textbf{Recommendation:}
  Derive the surface gradient analytically from a single noise evaluation, or
  reduce to a forward difference (3 evaluations instead of 5) — the wave normal
  is already an approximation. Alternatively cache the centre height and reuse it
  across both the displacement and gradient. Visual result is preserved within
  the existing approximation tolerance.

\item[\textbf{[SH9]}] \textbf{Caustics Computed When Intensity Is Zero}
  \hfill \severitybox{mediumcolor}{MEDIUM}

  \begin{tabular}{ll}
    \textbf{File:} & \texttt{shaders/water.frag} \\
    \textbf{Line:} & 416--474 \\
    \textbf{Category:} & Fragment ALU \\
  \end{tabular}

  \textbf{Description:}
  The full caustic computation (either the 27-cell \texttt{voronoi3d} loop or the
  Perlin Jacobian with its 5 FBM evaluations, see SH2) runs before
  \texttt{caustic} is scaled by \texttt{causticIntensity} at line 472. When a
  water material has caustics disabled (\texttt{causticIntensity == 0}), all of
  that work is discarded.

  \textbf{Recommendation:}
  Wrap the caustic block in \texttt{if (causticIntensity > 0.0)}. The value comes
  from the per-material SSBO and is uniform across a draw, so the branch is fully
  coherent and free where caustics are on. No visual change.

\end{enumerate}

\newpage

\subsection{Low Severity}

\begin{enumerate}[leftmargin=*]

\item[\textbf{[SH10]}] \textbf{Vegetation Background Mips Sampled Before Discard}
  \hfill \severitybox{lowcolor}{LOW}

  \begin{tabular}{ll}
    \textbf{File:} & \texttt{shaders/vegetation.frag} \\
    \textbf{Line:} & 53--55, 70 \\
    \textbf{Category:} & Texture Bandwidth \\
  \end{tabular}

  \textbf{Description:}
  Two \texttt{textureLod} background samples (\texttt{bgAlbedo},
  \texttt{bgNormEnc} at mip 5) are taken at lines 54--55, then the fragment may
  \texttt{discard} at line 70 when \texttt{weight < 0.3}. For the many
  transparent billboard texels that are discarded, the background fetches are
  wasted bandwidth.

  \textbf{Recommendation:}
  Move the two background \texttt{textureLod} fetches below the
  \texttt{weight < 0.3} discard test. They are only consumed at line 131, well
  after the discard, so reordering is safe and changes nothing visually. (Note:
  keep the opacity/albedo/normal fetches where derivatives are still valid, or
  switch background to explicit-LOD which is already the case.)

\item[\textbf{[SH11]}] \textbf{$O(n^2)$ Refraction Blur Tap Loop}
  \hfill \severitybox{lowcolor}{LOW}

  \begin{tabular}{ll}
    \textbf{File:} & \texttt{shaders/water.frag} \\
    \textbf{Line:} & 228--244 \\
    \textbf{Category:} & Texture Bandwidth \\
  \end{tabular}

  \textbf{Description:}
  When \texttt{enableBlur} is set, the refraction blur uses a nested
  \texttt{(2\,halfK+1)$^2$} loop of \texttt{sceneColorTex} samples. For
  \texttt{blurSamples = 5} that is 25 dependent texture reads per fragment, all
  with an unrolled loop. This is optional (guarded by \texttt{enableBlur}), hence
  low severity, but it dominates cost when enabled.

  \textbf{Recommendation:}
  Replace the box blur with a separable two-pass Gaussian, or reduce to a small
  fixed weighted kernel (e.g.\ 5-tap cross). Since the water refraction is
  already a soft effect, a separable approximation keeps the look while turning
  $n^2$ taps into $2n$.

\item[\textbf{[SH12]}] \textbf{Water Debug Mode 32 Recomputes Displacement}
  \hfill \severitybox{lowcolor}{LOW}

  \begin{tabular}{ll}
    \textbf{File:} & \texttt{shaders/water.frag} \\
    \textbf{Line:} & 481--509 \\
    \textbf{Category:} & Fragment ALU (debug path) \\
  \end{tabular}

  \textbf{Description:}
  Debug mode 32 recomputes \texttt{waterWaveDisplacement} (another full FBM
  evaluation) purely to visualize displacement, duplicating work already done by
  the TES and available in \texttt{fragDebug}. This is subsumed by SH7 once the
  debug ladder is specialization-gated, but is noted separately because it adds
  a heavy FBM call inside the debug path.

  \textbf{Recommendation:}
  Fold into the SH7 specialization-constant gating; when kept, prefer the
  interpolated \texttt{fragDebug} value the shader already receives.

\end{enumerate}

\newpage

\section{Summary Table}

\begin{center}
\begin{tabular}{llll}
\toprule
\textbf{ID} & \textbf{Shader} & \textbf{Severity} & \textbf{Nature} \\
\midrule
SH1  & water.frag              & HIGH   & Per-fragment \texttt{inverse(mat4)} \\
SH2  & water.frag / water\_noise & HIGH & Redundant FBM/caustic noise \\
SH3  & perlin.glsl / voronoi.glsl & HIGH & \texttt{sin()}-based hash \\
SH4  & water.frag              & MEDIUM & Repeated screenUV / depth samples \\
SH5  & main.frag               & MEDIUM & Unconditional env cubemap fetch \\
SH6  & triplanar.glsl          & MEDIUM & UV recompute + repeated sampling \\
SH7  & main.frag / water.frag  & MEDIUM & Debug ladder in release build \\
SH8  & water.tese              & MEDIUM & 5$\times$ displacement per vertex \\
SH9  & water.frag              & MEDIUM & Caustics run when intensity 0 \\
SH10 & vegetation.frag         & LOW    & BG mips fetched before discard \\
SH11 & water.frag              & LOW    & $O(n^2)$ blur tap loop \\
SH12 & water.frag              & LOW    & Debug 32 recomputes displacement \\
\bottomrule
\end{tabular}
\end{center}

\section{Overall Classification}

The shader suite is functionally rich and correct, but the two procedural-heavy
fragment shaders carry substantial redundant per-fragment cost. The highest-
impact wins require \textbf{no visual change whatsoever}:

\begin{itemize}[noitemsep]
  \item \textbf{SH1} (hoist the matrix inverse to a UBO) is a near-free CPU-side
        change removing $\sim$200 ops per water fragment.
  \item \textbf{SH3} (integer hash) accelerates \emph{all} procedural noise —
        water, terrain displacement and vegetation — by removing dozens of
        \texttt{sin} evaluations per noise octave.
  \item \textbf{SH2/SH8/SH9} attack the water noise budget, which is by far the
        most expensive per-fragment / per-vertex workload in the engine.
\end{itemize}

\medskip
\noindent\severitybox{mediumcolor}{OPTIMIZABLE}

\begin{quote}
  The shaders are validation-clean and feature-complete. Performance headroom is
  significant but entirely recoverable through invariant hoisting, hash
  replacement, coherent early-outs and specialization-constant gating of debug
  code. None of the recommended changes alter the rendered image.
\end{quote}

\subsection*{Remediation Priority}

\begin{enumerate}[noitemsep]
  \item[\textbf{P0}] SH1 (matrix inverse $\to$ UBO), SH3 (integer hash).
  \item[\textbf{P1}] SH2 (caustic noise reduction), SH9 (caustic early-out),
        SH5 (env fetch guard).
  \item[\textbf{P2}] SH4 (dedupe screenUV/depth), SH6 (triplanar UV reuse),
        SH8 (tese displacement), SH7 (debug specialization constants).
  \item[\textbf{P3}] SH10, SH11, SH12.
\end{enumerate}

\end{document}
